\makeabstract
    {
    Interleukin-6 licenses mast cell regranulation through CD71 upregulation after IgE-mediated activation in mice 
    }
    {
    Stephanie Park\textsuperscript{1,2}, Kenshiro Matsuda\textsuperscript{2,3}, Rizza U. Santos\textsuperscript{2}, Daniel F. Dwyer\textsuperscript{2,3}
    }
    {
    \textsuperscript{1} Amherst College, Amherst, MA \\
\textsuperscript{2} Division of Allergy and Clinical Immunology, Brigham and Women{\textquoteright}s Hospital, Boston, MA \\
\textsuperscript{3} Harvard Medical School, Boston, MA
    }
    {
    Mast cells (MCs) are immune cells that contribute to allergic inflammation by releasing mediators including histamine, tryptase, and leukotrienes upon activation. After immunoglobulin E (IgE)-mediated degranulation, MCs can recover and replenish their secretory granules through regranulation. We previously identified that CD71 expression in MCs was upregulated following IgE-driven activation. Therefore, this study investigated how CD71, the transferrin receptor for cellular iron uptake, is regulated during MC recovery after IgE activation.
    }
    {
    Bone marrow cells were cultured with 30 ng/mL stem cell factor (SCF) and 10 ng/mL interleukin (IL)-3 for 4{\textendash}6 weeks to generate bone marrow-derived cultured MCs (BMMCs). Cells were sensitized with 1 {\textmu}g/mL anti-dinitrophenyl-bovine serum albumin (DNP)-IgE overnight and activated with 0.2 {\textmu}g/mL DNP-BSA. Degranulation was assessed 1 hour after activation, and regranulation was assessed at 1 and 7 days by toluidine blue staining. CD71 expression during recovery was measured by flow cytometry with or without IL-6 treatment. Chromatin immunoprecipitation sequencing (ChIP-seq) datasets were analyzed to identify IL-6-associated transcription factor binding near the mouse CD71 locus. CRISPR-Cas9-mediated gene editing was used to disrupt CD71 in BMMCs, followed by assessment of viability and regranulation.
    }
    {
    BMMCs retained the ability to regranulate following IgE-mediated activation. Degranulation was evident 1 hour after activation and granules were restored by day 7. ChIP-seq analysis identified binding sites for multiple IL-6-associated transcription factors, including signal transducer and activator of transcription (STAT) 1, STAT3, STAT5A, STAT5B, JUN, ELK1, and CCAAT/enhancer-binding protein beta (C/EBP\ensuremath{\beta}), overlapping an active enhancer region near the mouse CD71 locus. CD71 expression increased following activation and gradually declined over time, but remained elevated through day 7 with IL-6 treatment. CRISPR-Cas9-mediated disruption of CD71 did not substantially affect BMMC viability but impaired restoration of granules after IgE-mediated activation compared with untransfected and off-target controls.
    }
    {
    These findings identify IL-6 as a critical regulator of MC regranulation following IgE-mediated activation by promoting CD71 expression. IL-6-associated transcriptional regulation may promote CD71 expression during recovery, potentially supporting iron-dependent processes required to restore MC granules. This pathway may help explain how MCs recover their secretory capacity following activation and remain capable of subsequent inflammatory responses.
    }

    \newpage
    